\scanpage{218}
\begin{center}- \mspage{4} -\end{center}

cation $\underline{P'^{q_o}}/(\textcircled{$m$}_{\underline A}\underline{P'^{q_o}})\longrightarrow\underline{P^{q_o}}/(\textcircled{$m$}_{\underline A}\underline{P^{q_o}})$ étant injective, il en résulte que le complexe
\diagram{\begin{tikzcd}[ampersand replacement=\&, column sep=small]
\cdots\arrow[r]\&\underline{P^{q_o-2}}\arrow[r]\&\underline{P^{q_o-1}}\arrow[r]\&\underline{P'^{q_o}}\arrow[r]\&0\arrow[r]\&\cdots
\end{tikzcd}}
vérifie les conditions du lemme. On peut donc supposer dans l'énoncé que $\underline{P^q}=0$ pour $\underline q>\underline q_{\underline o}$. Utilisons [(1), Ch.III, §.5] pour montrer que $\underline H$ est $\underline A$-plat. Le complexe $\underline P^\bullet$ est une résolution pour $\underline H$ par des $\underline A$-modules plats. Par consequence, les modules $\operatorname{Tor}^{\underline A}_{\bullet}(\underline A/\textcircled{$m$}_{\underline A},\underline H)$ sont les modules de cohomologie du complexe $\underline P^\circ\otimes_{\underline A}\underline A/\textcircled{$m$}_{\underline A}=\underline P^\circ/\textcircled{$m$}_{\underline A}\underline P^\circ$ et la conclusion résulte.

\uline{Corollaire}. \uline{Soit $\underline f:\underline X\longrightarrow\underline Y$ un morphisme propre et plat de variétés algébriques. Alors l'ensemble}
\[
\underline{\underline G_{\underline f}(\underline Y)=\{\underline y\in\underline Y\mid\underline X_{\underline y}\text{ soit de Gorenstein}\}}
\]
\uline{est ouvert dans $\underline Y$}.

\begin{flushright}
L'Institut Central de Mathématique\\
14, rue de l'Académie , Bucarest\\
ROUMANIE
\end{flushright}

\figure{\begin{tikzpicture}
\draw[line width=.8pt] (0,0)--(.35,0);
\node[anchor=west] at (.7,0) {latine};
\draw (0.16,-.65) circle (.18);
\node[anchor=west] at (.7,-.65) {gothique};
\draw (0.16,-1.3) ellipse (.17 and .2);
\draw (0.13,-1.3) ellipse (.17 and .2);
\node[anchor=west] at (.7,-1.3) {grec};
\draw[gray] (0.16,-2) ellipse (.17 and .2);
\node[anchor=west] at (.7,-2) {ronde};
\draw (0,-2.65) .. controls (.2,-2.75) and (2,-2.65) .. (2.6,-2.65) .. controls (3.2,-2.65) and (3.2,-2.3) .. (3.2,-2) -- (3.2,.3);
\draw (3.45,.5) .. controls (5,.5) and (8,.5) .. (9.5,.5) .. controls (9.7,.5) and (9.9,.5) .. (10,.4);
\end{tikzpicture}}

\uline{Bibliographie}

[1] Bourbaki, N. -- Algèbre commutative, Paris, Hermann \unsure{1961|1967}

[2] Grothendieck, A -- E.G.A., IV, 1965

[3] Samuel, P -- Seminaire d'Algèbre commutative, 1966/67.

[4] J. P. Serre -- Faisceaux algébriques cohérents, Ann. Math. 1955
\margin{217}
